%! Representing a Categorical Variable — Unit 1, Lesson 1.2 — Lesson Plan
% GRADUAL-RELEASE plan (warm-up / guided notes and practice / debrief / close and start homework).
% THERE IS NO GROUP ACTIVITY IN THIS COURSE — the guided notes carry the release: I do (two
% sections) -> we do (guided practice). THERE IS NO INDEPENDENT PRACTICE SET — the homework is
% the individual practice, and the period ends with students starting it. Teacher-facing. Every box is
% `skillbox` (breakable) with a \boxguard ahead of it; this course has no `fixedskillbox`,
% no tier boxes, and no exit ticket.
\documentclass[10pt]{article}
\usepackage{apstats-article}
\usepackage{apstats-boxes}

\newcommand{\UnitNumberName}{Unit 1: Exploring One-Variable Data \quad}
\newcommand{\LessonNumberName}{Lesson 1.2: Representing a Categorical Variable}

\begin{document}
\begin{center}
    {\Large\bfseries \CourseName \\
    {\small\bfseries \UnitNumberName \LessonNumberName}}
\end{center}

\begin{tcolorbox}[colback=sky, colframe=navy, boxrule=0.9pt, arc=2mm,
    left=2mm, right=2mm, top=1.5mm, bottom=1.5mm]
    \textbf{Primary Objective:} Students will build frequency and relative frequency tables for
    a categorical variable, read a bar graph of either one, and explain why two groups of
    different sizes can only be compared with relative frequencies.
    \\[2pt]
    \textbf{CED alignment:} Topics 1.3 \emph{and} 1.4 \quad$\cdot$\quad Big Idea: UNC
    (Uncertainty) \quad$\cdot$\quad Skills 2.A, 2.B, 2.D \quad$\cdot$\quad
    LO UNC-1.A--UNC-1.E \quad$\cdot$\quad EK UNC-1.A.1, 1.B.1, 1.B.2, 1.C.1--3, 1.D.1, 1.E.1
    \\[2pt]
    \textbf{Scope note:} this lesson absorbs \emph{two} CED topics (tables and graphs), per
    \texttt{COURSE\_RESTRUCTURE.md}. That is a lot for 55 minutes. Notes section~1 carries both
    tables, section~2 keeps the bar-graph mechanics to one pre-drawn chart, and section~2's second
    half --- the idea both topics share, \textbf{counts versus rates} --- gets the time. If the
    notes run long, trim the bar-graph discussion, never the crux that closes section~2.
    \\[2pt]
    \textbf{Lesson model:} \emph{Gradual release --- I do, we do, you do --- all of it inside
    the guided notes.} There is no group activity. The notes deliver both tables, the bar graph,
    and the comparison rule directly, in two sections; the guided practice is worked together
    with students holding the pen; \textbf{the homework is the individual practice}, and students
    start it in class while you circulate. The whole-class debrief is \emph{spoken} --- there is
    no transfer set and no reflection box on the page.
\end{tcolorbox}

\renewcommand{\arraystretch}{1.4}
\boxguard
\begin{skillbox}[Learning Targets \& Key Understandings]{goldbox}
\begin{tabularx}{\linewidth}{@{}X|X@{}}
\textbf{Learning targets (student language)} & \textbf{Key understandings (the ``why'')} \\[2pt]
\hline
\vspace{-6pt}
\begin{itemize}[leftmargin=1.1em, nosep, after=\vspace{2pt}]
  \item I can turn a table of counts into a table of shares, and say what each column must
        add to.
  \item I can read a bar graph and say which number each bar shows.
  \item I can explain why comparing different-sized groups by counts is unfair --- and when a
        count is still the right number.
  \item I can decide whether a ``most people'' claim is supported by the numbers.
\end{itemize}
&
\vspace{-6pt}
\begin{itemize}[leftmargin=1.1em, nosep, after=\vspace{2pt}]
  \item A frequency table gives counts; a relative frequency table gives proportions
        (EK UNC-1.A.1). Percents, proportions, and rates carry identical information
        (EK UNC-1.B.1).
  \item \textbf{Target misconception:} that the group with the larger count is ``more'' on that
        category. South has \emph{more walkers} (150 vs.\ 80) and a \emph{lower walking rate}
        (15\% vs.\ 20\%). Both are true; the bigger group wins on counts before anyone counts
        anything (EK UNC-1.E.1).
  \item \textbf{Second misconception:} that the largest category is a majority. Bus is 45\% ---
        a plurality, not ``most.''
  \item \textbf{Third misconception:} percent of the wrong total --- dividing by the two schools
        combined rather than by the one school being described.
\end{itemize}
\end{tabularx}
\end{skillbox}

\boxguard
\begin{skillbox}[{Vocabulary, Concepts, \& Theorems --- taught directly in the guided notes}]{greenbox}
\begin{minipage}[t]{0.48\linewidth}
    \begin{itemize}
        \item \textbf{Frequency table} --- the number of cases in each category (EK UNC-1.A.1)
        \item \textbf{Relative frequency table} --- the share of the group in each category;
              entries always sum to 1
        \item \textbf{Proportion / percent / rate} --- three presentations of one number
              (EK UNC-1.B.1)
    \end{itemize}
\end{minipage}
\hfill
\begin{minipage}[t]{0.48\linewidth}
    \begin{itemize}
        \item \textbf{Bar graph} --- bar height is the count or the relative frequency for a
              category; bars are separated by gaps (EK UNC-1.C.1, UNC-1.C.2)
        \item \textbf{The comparison rule} --- groups of different sizes are compared with
              relative frequencies, never raw counts (EK UNC-1.E.1)
        \item \textbf{Plurality vs.\ majority} --- largest share vs.\ more than half
        \item \textbf{Carried forward from Lesson 1.1} --- categorical vs.\ quantitative
    \end{itemize}
\end{minipage}
\end{skillbox}

% A tabularx never splits, so guard generously ahead of it.
\boxguard[30]
\begin{skillbox}[Lesson at a Glance --- \MeetingLength]{sky}
\renewcommand{\arraystretch}{1.25}
\begin{tabularx}{\linewidth}{@{}l c X X@{}}
\textbf{Phase} & \textbf{Min} & \textbf{Students} & \textbf{Teacher} \\
\hline
Warm-Up & 5 & Percents, and two teams with different game counts & Debrief item~3 aloud; leave their ``why the win totals did not decide it'' sentence on the board all period \\
Guided Notes \& Practice & 34 & Fill the vocabulary box and notes 1--2; then \emph{Two Shifts} together, pen in their hand & \textbf{I do} notes 1--2 (20) $\to$ \textbf{we do} the Guided Practice box (14) \\
Debrief & 8 & Pens down, papers up --- answer aloud, nothing to fill in & Put both true statements about walking side by side, then take the plurality and wrong-total traps \\
Close \& Assign & 8 & \textbf{Start the homework in class}, alone and silent & Announce the paper homework, circulate the first minutes for your formative read, preview Lesson~1.3 \\
\end{tabularx}
\end{skillbox}

\boxguard[20]
\begin{skillbox}[Warm-Up --- Activate Prior Knowledge (5 min)]{sky}
\begin{minipage}[t]{0.48\linewidth}
    \textbf{The three items and what each seeds}
    \begin{itemize}
        \item \textbf{1.} Classify species and weight. \textbf{Seeds:} the Lesson~1.1 spiral,
              and it establishes that today's variable is the categorical one.
        \item \textbf{2.} $80/400$ and $150/1000$ as percents. \textbf{Seeds:} the exact
              arithmetic the crux at the end of notes section~2 needs, done before the crux so
              it is not the obstacle.
        \item \textbf{3.} Two teams, different numbers of games. \textbf{This is the whole
              lesson in prior-knowledge terms} --- every student already knows you compare
              records by rate.
    \end{itemize}
\end{minipage}
\hfill
\begin{minipage}[t]{0.48\linewidth}
    \textbf{Running it}
    \begin{itemize}
        \item Item~3's numbers are the same numbers the crux in section~2 uses. \textbf{Debrief
              it aloud} and press the second half: \emph{why} the raw win totals did not decide
              it. Write their phrasing on the board verbatim and leave it up.
        \item You will point at that sentence twice more --- when the work block in section~2
              lands, and again in the closing debrief.
        \item If students are slow on item~2's percents, spend the extra minute here rather than
              inside the notes.
    \end{itemize}
\end{minipage}
\end{skillbox}

\boxguard
\begin{skillbox}[Guided Notes \& Practice --- I do, then we do (34 min)]{sky}
\textbf{This is the whole lesson.} There is no group activity, and there is no independent
practice set on this page: \textbf{the homework is the individual practice}, started in class.
\begin{multicols}{2}
\textbf{I do (about 20 min).} Two sections, each in two moves; students fill the blanks as each
idea is named, and the vocabulary box is filled \emph{as you go}, never in advance.

\emph{Section 1} builds both tables off one column and lands the two totals as \textbf{two
different promises} --- 400 students, and 1. Take the plurality trap right there: bus at 45\% is
the largest share and is not ``most.'' Its second half, \emph{Proportion, Percent, Rate}, is one
line --- $80/400 = 0.20 = 20\%$ --- plus the AP habit of attaching context. Keep that half under
two minutes. \emph{Section 2} opens with the only graph block: bar $=$ category, height $=$ count
\emph{or} relative frequency, gaps because categories are not a number line. Use the \emph{Read
it} fills to make them take numbers off the chart rather than watch you do it. \textbf{Section
2's second half is the lesson.} Point at the posted side-by-side chart before computing anything:
South is taller almost everywhere, and that was guaranteed. Then work the \texttt{work} block on
the board and put both true statements up together. Close on the counterweight --- ``how many
students walk in this district'' wants a count.

\textbf{We do (about 14 min).} The Guided Practice box, \emph{Two Shifts}. Students hold the pen;
you hold the questions: \emph{``Which shift sold more? Which sells them more often? Are those the
same question?''} \quad \emph{``Busier --- is that about how many, or about what share?''} Part
(c) is the counterweight again and is the part to protect when time is short.

\columnbreak
\textbf{The pen must actually change hands.} Guided Practice is the only place today where you
can watch a student decide, so do not narrate it. Put part (a) on them cold, take a hand count on
\emph{which shift sells large drinks at the higher rate} before anyone speaks, and make the room
defend both numbers.

Circulate during Guided Practice and demand the reason, not the letter:
\begin{itemize}[leftmargin=1.1em, nosep]
  \item \textbf{Part (a), the crux:} ``Go back to the warm-up. Team B won more games. Was Team B
        better?''
  \item \textbf{Part (b):} ``Which kind of number did you write the sentence with --- and why
        that one?''
  \item \textbf{Part (c):} ``Busier --- how many, or what share? Point at the column.''
\end{itemize}
While circulating, pick the papers to put up in the debrief --- one that answered part (a) with
the count alone, one that justified the rate correctly. \textbf{Your formative read comes twice}:
here, and again in the last eight minutes as students start the homework.
\end{multicols}
\end{skillbox}

\boxguard
\begin{skillbox}[Debrief --- whole class (8 min)]{sky}
Pens down, papers up. \textbf{This debrief is spoken} --- there is no box at the end of the
notes to fill in, so it lives entirely on the board and in the room.
\begin{multicols}{2}
\begin{enumerate}[leftmargin=1.3em, nosep, itemsep=3pt]
  \item \textbf{Put both true statements side by side on the board:} \emph{South has more
        walkers} (150 vs.\ 80) and \emph{North has the higher walking rate} (20\% vs.\ 15\%).
        Make a student say which one answers ``which school walks more,'' and hear the answer
        \emph{it depends on the question}.
  \item Circle the warm-up sentence still on the board from item~3 and write the rule next to
        it: comparing different-sized groups requires relative frequencies (EK UNC-1.E.1).
  \item Put up the two Guided Practice papers you picked and name what the correct one did.
  \item Take the plurality trap next, off notes~1 --- 45\% is the largest share, and 55\% of
        North does something else. Land it as a reading habit, not a vocabulary item.
  \item Take the wrong-total trap last: \emph{percent of what?} Point at the two-schools table in
        notes~2 and ask what dividing by 1,400 would have computed.
\end{enumerate}
\columnbreak
\textbf{Two things to demand aloud.} First, \textbf{make someone name a question in this lesson
that a count answers correctly} --- the Total column in \emph{Two Shifts}, or ``how many students
in the district walk.'' Students who over-learn ``always use rates'' need this, and without it
the rule becomes a reflex instead of a decision. Second, \textbf{make them state the comparison
rule and then produce a pair of groups that were not on any page today} where counts would
mislead; that is where transfer shows up now that there is no transfer set on paper. Hold the
room to both rather than letting students race ahead to the homework --- they will start it in
eight minutes.

\textbf{Connect back to Lesson~1.1 if time allows:} every number on today's pages came from
\emph{counting} a categorical column. That is exactly what yesterday said categorical variables
are summarized by --- and it is why none of today's percents make the variable quantitative.

\textbf{If the debrief runs short}, cut the Lesson~1.1 callback and the wrong-total trap ---
item~1 is the only one that cannot be cut. \textbf{Do not let the debrief eat the last eight
minutes} --- the homework start is not optional padding; it is the only individual practice in
the lesson and your second formative read.
\end{multicols}
\end{skillbox}

\boxguard
\begin{skillbox}[AP Practice --- assigned, not scored]{goldbox}
Four AP-style multiple-choice items on a city recycling audit --- a rate computation, the crux
stated as four competing claims, the structural must-sum-to-1 fact, and a display-choice item
(EK UNC-1.E.1) --- then one multi-part free-response set on a 9th-versus-12th-grade activity
survey: identify and classify, build a relative frequency table with the arithmetic shown,
correct a newspaper's count-based comparison, and describe the largest difference in context.
\textbf{Part (a) is the spiral} --- it sends students back to Lesson 1.1 to classify the variable
before they are allowed to summarize it, which is the habit that keeps them from averaging a
categorical column later.

\textbf{Not scored --- the cover's score column reads \textbf{NA} for this page.} \textbf{This
page now carries the transfer:} the recycling audit and the activity survey appear nowhere else
in the packet, so treat what comes back in Lesson~1.3 as your real evidence that the comparison
rule transferred. Go over it at the start of Lesson~1.3 and hold the AP standard out loud:
\emph{in context or it does not count}.
\end{skillbox}

\boxguard
\begin{skillbox}[Homework --- scored, due next class]{goldbox}
Five items in fresh contexts: \textbf{1} a club vote --- percent plus the ``most'' check;
\textbf{2} the \emph{crux again} in a medical context, with both rates required; \textbf{3} the
same data plotted as counts and as proportions --- what changed and what did not (EK UNC-1.C.1);
\textbf{4} four percents adding to 95\%, a structural check; \textbf{5} an AP-style
multiple-choice item with a required justification. The packet closes with the Lesson~1.3
preview.

\textbf{Start it in the last eight minutes of the period --- this is the lesson's individual
practice, not homework tacked on the end.} \textbf{Scoring:} 10 points --- 2 per item. \textbf{Item~5 carries the check}: its distractors are the two halves of the crux stated
separately, so a student who has only half the idea will pick (A) or (B). Item~2 is the crux
transplanted into a medical context and is worth the most partial credit --- insist on both rates
written out, not just the word ``Ridgeview.''

\textbf{Paper homework is the default for this course} --- assign this page unless you have
decided otherwise. \textbf{DeltaMath override (occasional):} on the days you assign DeltaMath
instead, it replaces this page, you strike through the score cell on the cover, and this page
stays in the packet as extra practice. Never assign both.
\end{skillbox}

\boxguard
\begin{skillbox}[Watch For (while circulating)]{redbox}
\begin{itemize}[leftmargin=1.2em, nosep]
  \item \textbf{Counts compared across different-sized groups} (notes~2, GP~(a), HW~2 and~5,
        AP~MC2 --- the target misconception). Probe: \emph{``Team B won more games. Was Team B
        better?''}
  \item \textbf{Plurality read as majority} (notes~1, HW~1). Probe: \emph{``45\% chose the
        bus. What did the other 55\% do?''}
  \item \textbf{Percent taken of the wrong total} (notes~2) --- dividing by 1,400 (both schools)
        or by the category rather than the group. Probe: \emph{``Percent \emph{of what}? Point at
        the number you divided by.''}
  \item \textbf{Relative frequencies not summing to 1} (notes~1, HW~4, AP~MC3): students treat
        the column as four unrelated numbers. Probe: \emph{``Add your column. What should it come
        to, and why must it?''}
  \item \textbf{Rate applied where a count is wanted} (GP~(c)): over-correction after
        the crux lands. Probe: \emph{``Busier --- how many, or what share?''}
  \item \textbf{Answers with no context:} ``20\%'' alone. Probe: \emph{``20\% of whom, doing
        what?''} Hold this from day one; it is the AP scoring line.
\end{itemize}
Cold-call prompts: \emph{``Give me one number from this table that is a count and one that is a
rate.''} \quad \emph{``What would make this chart a fair comparison?''}
\end{skillbox}

\boxguard
\begin{skillbox}[Close \& Assign (8 min)]{goldbox}
\textbf{Students start the homework here, in class, alone and silent --- this is the individual
practice, and the first minutes of it are your best formative read.} Hand the launch first ---
five items on paper, scored, due next class --- and point out that the AP Practice page is theirs
to work and bring to review. Circulate: item~2, the two hospitals, is the column that tells you
whether the crux at the end of section~2 landed. Sort what you see into three piles as you walk
--- chose the rate \emph{and} named why (ready); chose the rate with no reason (ready, needs the
language); chose the count (the misconception survived) --- and open Lesson~1.3 by reteaching if
that third pile ran past a third of the room. Then close on what changed:
\emph{``You walked in able to read a table. You walk out able to catch a chart that is
technically accurate and still misleading --- which is most of the charts you will ever be
shown.''} \textbf{Preview:} Lesson~1.3, \emph{Representing a Quantitative Variable with Graphs}
--- counting worked today because there were four category names. When the values are numbers
that barely repeat, counting stops working and we need different displays. Note that
\emph{segmented bar charts} and \emph{pie charts} are named in the CED (EK UNC-1.C.3) but are
deliberately \textbf{not} built here; they return in Unit 2 with two-way tables, where they
actually earn their place.
\end{skillbox}

% ── Teacher notes — one per component, in packet order (four: there is no activity) ──
\begin{teachernote}[Warm-Up]
Item~3 is the one that matters --- it is the entire lesson in language every student already
has, and its numbers are deliberately \emph{not} the ones the crux in section~2 uses, so the transfer is
real. Give three minutes and make sure the second half gets answered out loud: \emph{why} Team
B's larger win total did not decide it. \textbf{Write that sentence on the board verbatim and
leave it up all period}; you will circle it in the debrief. Item~2 exists so the arithmetic is
not the obstacle later --- if students are slow on percents, spend the extra minute here rather
than inside the notes. Item~1 is a 60-second spiral to Lesson 1.1.
\end{teachernote}

\begin{teachernote}[Guided Notes \& Practice]
\textbf{This one document is the lesson --- pace it 20 / 14, then 8 for the debrief, then 8 for
the homework start.}
\textbf{This lesson covers two CED topics, so the I-do is the tightest of the unit.} Section~1
is the tables and must move: the two-table build in about six minutes, the
proportion-percent-rate half in two. The bar-graph half that opens section~2 is the only stretch
of pure display instruction --- keep it under four minutes and keep pointing at the printed
chart; the \emph{Read it} fills are what stop it becoming a lecture.

\textbf{Spend the time on the second half of section~2}, which carries the misconception. Before
you compute anything, point at the district's side-by-side chart and ask whether South could have
lost any category. Then work the \texttt{work} block live and put both true statements on the
board at once --- \emph{South has more walkers} and \emph{North has the higher rate}. Do not
resolve the tension too fast; the sentence to land is that they answer two different questions.

\textbf{Do not skip the counterweight} that closes section~2 (``how many students in this
district walk?''). Without it students leave with ``always use rates,'' which is wrong, and they
will over-apply it on the AP exam.

\textbf{Guided Practice is where the pen changes hands.} \emph{Two Shifts} reverses the numbers:
morning sells more large drinks (84 vs.\ 32) while evening sells them at the higher rate (40\%
vs.\ 35\%). Part (c) is the counterweight in its cleanest form and is the part to protect when
you are running behind.

\textbf{There is no independent practice set on this page, and that is deliberate --- the
homework is the individual practice.} Guided Practice is the last thing on the page; the period
ends with students starting the homework in front of you, which is where your formative read
comes from. Item~2, the two hospitals, is the item to watch: sort as you circulate into
chose-the-rate-with-a-reason (ready), chose-the-rate-with-no-reason (ready, needs the language),
and chose-the-count (the misconception survived). If the third pile is more than a third of the
room, \textbf{open Lesson~1.3 by reteaching the two-schools table}, and say so plainly.

\textbf{Debrief (8 min), spoken.} Both-statements-on-the-board first; it is the one move that
cannot be cut. Making a student name a question that a \emph{count} answers correctly is the
second-most valuable minute in the debrief. \textbf{Do not let the debrief eat the last eight
minutes} --- the homework start is not optional padding.
\end{teachernote}

\begin{teachernote}[AP Practice]
\textbf{This page is not scored} --- the graded practice for this lesson is the homework that
follows it. Work it as AP-format reps and go over it at the start of Lesson~1.3. \textbf{This
page now carries the transfer:} the recycling audit appears nowhere else in the packet, so the
four multiple-choice items are the first time students meet this data set.

\textbf{MC item~2 is the one worth reading closely}: options (A) and (B) are each half of the
truth, and students who pick either have not held both facts at once. On the free response,
part (b) is where students first have to \emph{show} relative frequency arithmetic rather than
state a result --- hold the line on writing the division, since the AP exam awards the setup.
Part (c) is the highest-value part: it asks them to correct someone else's statistic, which is
the form this idea takes on the exam and in every news story they will read. A part (b) with
three decimals and no context earns nothing; a part (d) that says ``study hall is different''
without the two percents is incomplete.
\end{teachernote}

\begin{teachernote}[Homework]
\textbf{Scored, 10 points (2 per item), due next class.} Launch it in the last eight minutes so
students hit item~1 while you can still answer --- \textbf{this page is the lesson's individual
practice}, and the first minutes of it are your second formative read.

\textbf{Item~2 is the one to read closely.} The two hospitals are the crux transplanted into a
context with much larger and much more lopsided groups (340 of 8,500 against 60 of 1,000), and
the item is worth the most partial credit --- insist on both rates written out, not just the
name ``Ridgeview.'' Item~4 is the structural check and takes ten seconds if students have
internalised that relative frequencies must sum to 1; if it is taking longer than that, the
section~1 fills did not land. \textbf{When you score the page, sort item~5 into three piles} ---
chose (C) with both rates computed (ready); chose (C) by elimination (ready, needs the
arithmetic); chose (A) or (B) (the misconception survived). If more than a third land in the
last pile, reopen the two-schools table in the Lesson~1.3 warm-up debrief rather than letting it
ride. Hold the context line while scoring: a percent with no ``of whom, doing what'' earns
nothing on the AP exam and should earn nothing here.

\textbf{Paper is the default.} This page is what gets assigned unless you decide otherwise on a
given day. On the occasions you assign DeltaMath in its place, strike the score cell on the cover
and tell students this page is now optional practice. Do not assign both.
\end{teachernote}
\end{document}
